% !TeX root = compte-rendu.tex
% !TeX spellcheck = fr


\section{Parlement Européen}

L'Europe essaie comme toutes les démocraties, d'appliquer le principe "une personne, une voix". Cependant, comme le disent Werner Kirsch, Moshé Machover, Wojciech Słomczyński et Karol Życzkowski \cite{penrose62}, le système est en 2 étapes : les citoyens élisent leurs élus puis ceux-ci votent au Parlement Européen. Il faut donc aussi que leurs représentants aient le même pouvoir.


\subsection{Équité entre pays membres}

Actuellement, le nombre de sièges attribués à chaque pays résulte des traités signés entre les pays. Il n'existe pas de formule précisant la répartition des sièges.

Pour évaluer le pouvoir d'un pays, imaginons que tout les eurodéputés de ce pays votent toujours ensemble. Cela revient à considérer que la possibilité des pays forment des coalitions dans leur intérêt propre et non dans celui de l'Union Européenne. \par
Chaque pays est un votant dont le poids est égal au nombre d'eurodéputés qu'il a. Il suffit d'exécuter les algorithmes décrits dans la \hyperlink{section.2}{partie 2}.

\begin{figure}[H]
	\begin{center}
		\begin{tabular}{|l|r|r|r|}
			\hline \diagbox[dir=NE]{~~~~~~~~~~~~}{~} & Nombre de sièges &
			\begin{tabular}{@{}c@{}}
				Pourcentage \\ de sièges
			\end{tabular} &
			\begin{tabular}{@{}c@{}}
				Indice de \\ Banzhaf normalisé
			\end{tabular} \\
			\hline Allemagne & 96 & 13,35\% & 14,32\% \\
			\hline France & 81 & 11,27\% & 11,61\% \\
			\hline Italie & 76 & 10,57\% & 10,81\% \\
			\hline Espagne & 60 & 8,34\% & 8,33\% \\
			\hline Pologne & 53 & 7,37\% & 7,33\% \\
			\hline Roumanie & 33 & 4,59\% & 4,47\% \\
			\hline Pays-Bas & 31 & 4,31\% & 4,20\% \\
			\hline Belgique & 22 & 3,06\% & 2,97\% \\
			\hline République tchèque & 21 & 2,92\% & 2,83\% \\
			\hline Grèce & 21 & 2,92\% & 2,83\% \\
			\hline Hongrie & 21 & 2,92\% & 2,83\% \\
			\hline Portugal & 21 & 2,92\% & 2,83\% \\
			\hline Suède & 21 & 2,92\% & 2,83\% \\
			\hline Autriche & 20 & 2,78\% & 2,70\% \\
			\hline Bulgarie & 17 & 2,36\% & 2,29\% \\
			\hline Danemark & 15 & 2,09\% & 2,02\% \\
			\hline Finlande & 15 & 2,09\% & 2,02\% \\
			\hline Slovaquie & 15 & 2,09\% & 2,02\% \\
			\hline Irlande & 14 & 1,95\% & 1,89\% \\
			\hline Croatie & 12 & 1,67\% & 1,61\% \\
			\hline Lituanie & 11 & 1,53\% & 1,48\% \\
			\hline Lettonie & 9 & 1,25\% & 1,21\% \\
			\hline Slovénie & 9 & 1,25\% & 1,21\% \\
			\hline Estonie & 7 & 0,97\% & 0,94\% \\
			\hline Chypre & 6 & 0,83\% & 0,81\% \\
			\hline Luxembourg & 6 & 0,83\% & 0,81\% \\
			\hline Malte & 6 & 0,83\% & 0,81\% \\
			\hline Total & 720 & & \\
			\hline
		\end{tabular}
	\end{center}
\end{figure}


\begin{figure}[H]
	\includegraphics[width=\textwidth]{diagramme_pouvoir_UE_pays.png}
	\caption{diagramme en camembert représentant les sièges et l'indice de pouvoir de Banzhaf que détient chaque pays au Parlement Européen}
\end{figure}

\begin{figure}[H]
	\includegraphics[width=\textwidth]{graph_écart_pourvoir_représentation_pays.png}
	\caption{diagramme montrant pour chaque pays l'écart relatif $\nicefrac{(pouvoir - repr\acute{e}sentation)}{repr\acute{e}sentation}$ où la représentation est la pourcentage de la population européenne du pays}
\end{figure}

Dans le système actuel du Parlement Européen, les trois grands pays l'Allemagne, la France et l'Italie sont surreprésentés. L'Allemagne a un écart relatif le plus grand qui vaut 7,29\%.
Les autres pays ont un écart relatif d'environ 3\%.


\subsection{Pouvoir des différents groupes}

Les membres non inscrits n'ont aucune raison de voter ensemble et sont situés à des endroits très différents dans le spectre droite-gauche. Ils ne peuvent donc pas être considérés comme un groupe.
Nous ne les considérerons pas dans notre étude.

\begin{figure}[H]
	\begin{center}
		\begin{tabular}{|l|r|r|r|r|}
			\hline \diagbox[dir=NE]{~~~~~~~}{~} & Nombre de sièges &
			\begin{tabular}{@{}c@{}}
				Pourcentage \\ de sièges
			\end{tabular} &
			\begin{tabular}{@{}c@{}}
				Indice de \\ Banzhaf
			\end{tabular} &
			\begin{tabular}{@{}c@{}}
				Indice de \\ Banzhaf réaliste
			\end{tabular} \\
			\hline GUE/NGL & 46 & 6,67\% & 5,65\% & 0\% \\
			\hline S\&D & 136 & 19,71\% & 18,55\% & 5,26\% \\
			\hline Verts/ALE & 53 & 7,68\% & 7,26\% & 5,26\% \\
			\hline RE & 75 & 10,87\% & 8,87\% & 5,26\% \\
			\hline PPE & 188 & 27,25\% & 31,45\% & 52,63\% \\
			\hline CRE & 80 & 11,59\% & 10,48\% & 15,79\% \\
			\hline PfE & 86 & 12,46\% & 13,71\% & 15,79\% \\
			\hline ENS & 26 & 3,77\% & 4,03\% & 0\% \\
			\hline Total & 690 & & & \\
			\hline
	\end{tabular}
	\end{center}
\end{figure}

\begin{figure}[H]
	\includegraphics[width=\textwidth]{diagramme_pouvoir_UE_groupe.png}
	\caption{diagramme en camembert représentant les sièges, l'indice de pouvoir de Banzhaf et l'indice de Banzhaf réaliste que détient chaque groupe au Parlement Européen}
\end{figure}

Le parti populaire européen joue un rôle central dans le Parlement Européen simplement grâce au nombre de sièges qu'il détient.

L'indice de Banzhaf réaliste accentue encore la différence de pouvoir entre le PPE et les autres groupe. Ni les groupes à sa gauche (310 membres) ni les groupes à sa droite (192 membres) n'ont la majorité. Les coalitions gagnantes contiennent donc nécessairement le PPE. 
Au contraire, le S\&D, bien que second groupe en terme de nombre de membres, est le deuxième groupe le plus à gauche. Les groupe des Verts/ALE aux ENS peuvent facilement former des coalitions sans le S\&D. Son indice de Banzhaf réaliste est ainsi faible.
De même, les groupes GUE/NGL et ENS situés aux extrémités du Parlement Européen ont un indice de Banzhaf nul puisque toute coalition gagnante réaliste où ils sont, contient aussi le PPE ce qui rend ces deux groupes non décisifs.

\begin{figure}[H]
	\includegraphics[width=\textwidth]{graph_écart_pourvoir_représentation_groupe.png}
	\caption{diagramme montrant pour chaque groupe l'écart relatif $\nicefrac{(pouvoir - repr\acute{e}sentation)}{repr\acute{e}sentation}$}
\end{figure}

\begin{figure}[H]
	\includegraphics[width=\textwidth]{graph_écart_pourvoir_réaliste_représentation_groupe.png}
	\caption{diagramme montrant pour chaque pays l'écart relatif $\nicefrac{(pouvoir~r\acute{e}aliste - repr\acute{e}sentation)}{repr\acute{e}sentation}$}
\end{figure}